% TOP_PROBLEM: {"id": 3169, "title": "Geodesic cycles and Tutte's Theorem", "category_id": 3, "category": {"id": 3, "name": "graph_theory", "display_name": "Graph Theory", "description": "Problems involving graphs, networks, and their properties.", "slug": "graph-theory", "order_index": 3, "created_at": "2026-07-31T15:26:25.671Z"}, "status": "open", "problem_number": "OPG-500", "set_id": 12}
% FIRST_POSTED: 2026-10-01
% ENTRY_KIND: statement_only
% UnsolvedMath ID 3169; source catalogue code OPG-500
% !TeX program = lualatex
% Definition and sources only; this file is not an LLM solution attempt.
\documentclass[11pt,a4paper]{article}
\usepackage[margin=25mm]{geometry}
\usepackage{fontspec}
\setmainfont{lmroman10-regular.otf}[BoldFont=lmroman10-bold.otf,
 ItalicFont=lmroman10-italic.otf,BoldItalicFont=lmroman10-bolditalic.otf]
\usepackage{amsmath,amssymb,mathtools}
\usepackage{unicode-math}
\setmathfont{latinmodern-math.otf}
\AtBeginDocument{\let\setminus\smallsetminus}
\usepackage{xurl,hyperref}
\hypersetup{colorlinks=true,urlcolor=blue}
\setcounter{secnumdepth}{0}
\setlength{\parindent}{0pt}
\setlength{\parskip}{5pt}
\setlength{\emergencystretch}{3em}
\begin{document}
\section*{Geodesic cycles and Tutte's Theorem}\label{3169}
{\small UnsolvedMath ID 3169 · OPG-500 · Graph Theory · OpenGarden}
\subsection{Definitions and mathematical statement}
Let $G$ be a finite simple three-vertex-connected graph. This means $|V(G)|\ge4$ and removing fewer than three vertices leaves a connected graph. Give each edge $e$ a strictly positive real length $\ell(e)$. For vertices $u,v$, the distance $d_G(u,v)$ is the minimum total edge length of a path joining them.

A cycle $C$ is geodesic when its intrinsic metric agrees with the graph metric on its vertices:
\[
       d_C(u,v)=d_G(u,v)\qquad(u,v\in V(C)),
\]
where $d_C(u,v)$ is the smaller of the lengths of the two $u$--$v$ arcs of $C$. Equal-length shortest paths are allowed; the definition does not require uniqueness or that every shortest path stay inside $C$.

A peripheral cycle is an induced cycle whose deletion leaves the graph connected. Induced means that no edge joins nonconsecutive vertices of the cycle, and deletion here means removing all vertices of $C$ and their incident edges. The conjecture asks whether every $G$ with every positive edge-length assignment has a cycle that is both peripheral and geodesic.

The cycle may depend on all the chosen edge lengths. Merely choosing a shortest cycle does not guarantee the peripheral property, and merely finding an induced nonseparating cycle does not guarantee the distance equalities. The target requires both properties of one cycle, with the distance test made against all paths in the original weighted graph.
\subsection{Short English statement}
Does every positively weighted three-connected graph contain an induced nonseparating cycle that preserves all distances between its vertices?
\subsection{Sources}
\begin{itemize}
\item[S1] ULAM.AI and the UnsolvedMath Contributors, supplied problems.json, record 3169 (OPG-500), OpenGarden. Catalogue identity and source transcription; CC BY 4.0.
\url{https://huggingface.co/datasets/ulamai/UnsolvedMath}.
\item[S2] Open Problem Garden, Geodesic cycles and Tutte's Theorem. Original problem and discussion.
\url{http://www.openproblemgarden.org/op/geodesic_cycles_and_tuttes_theorem}.
\end{itemize}
\end{document}
